\section*{Chapitre \ref{ch:g5:brain-in-command} --- Le cerveau aux commandes}

\begin{solution}{exo:g5:brain-in-command:1}
Le \omterm{def:g5:brain-in-command:system}{cerveau}~: il reçoit tous les rapports, prend toutes les décisions,
donne tous les ordres. La \omterm{def:g5:brain-in-command:system}{moelle épinière}~: le gros câble qui descend
du \omterm{def:g5:brain-in-command:system}{cerveau} à l'intérieur de la \omterm{def:g4:classifying-animals:vertebrate}{colonne vertébrale}. Les \omterm{def:g5:brain-in-command:system}{nerfs}~: les
fils ramifiés qui portent les messages entre le \omterm{def:g5:brain-in-command:system}{cerveau} ou la moelle
et chaque organe.
\end{solution}

\begin{solution}{exo:g5:brain-in-command:2}
Les \omterm{def:g1:the-five-senses:sense}{organes des sens} signalent au \omterm{def:g5:brain-in-command:system}{cerveau}~; le \omterm{def:g5:brain-in-command:system}{cerveau} décide et
compose des ordres~; les \omterm{def:g5:brain-in-command:system}{nerfs} portent les ordres aux \omterm{def:g3:bones-and-muscles:muscle}{muscles}, qui
se contractent. Les \omterm{def:g1:the-five-senses:sense}{organes des sens} signalent, le \omterm{def:g5:brain-in-command:system}{cerveau} décide,
les \omterm{def:g3:bones-and-muscles:muscle}{muscles} obéissent.
\end{solution}

\begin{solution}{exo:g5:brain-in-command:3}
Le \omterm{def:g5:brain-in-command:system}{cerveau} dans le \omterm{def:g3:bones-and-muscles:skeleton}{crâne}~; la \omterm{def:g5:brain-in-command:system}{moelle épinière} à l'intérieur de la
colonne de petits \omterm{def:g3:bones-and-muscles:skeleton}{os} empilés.
\end{solution}

\begin{solution}{exo:g5:brain-in-command:4}
Le délai entre un signal et la réponse du corps --- environ un
cinquième de seconde pour une tâche simple chez une personne
entraînée, un peu plus chez les enfants.
\end{solution}

\begin{solution}{exo:g5:brain-in-command:5}
La \omterm{def:g1:the-five-senses:sense}{peau} du pied a signalé (toucher et douleur), le \omterm{def:g5:brain-in-command:system}{système nerveux} a
décidé le retrait, et les \omterm{def:g3:bones-and-muscles:muscle}{muscles} de la \omterm{def:g1:my-body:parts}{jambe} ont obéi en retirant
le pied d'un coup. (Un retrait aussi rapide est en partie traité par
la \omterm{def:g5:brain-in-command:system}{moelle épinière} elle-même --- le guichet des urgences du centre
de commande.)
\end{solution}

\begin{solution}{exo:g5:brain-in-command:6}
Il entretient la \omterm{def:g4:breathing:movements}{respiration} jour et nuit, et règle le rythme du
\omterm{def:g4:heart-and-blood:heart}{cœur} sur les besoins du corps (aussi~: garder les souvenirs, rêver).
\end{solution}

\begin{solution}{exo:g5:brain-in-command:7}
Le \omterm{prop:g1:healthy-habits:needs}{sommeil} --- son temps de classement et de réparation --- et une
grande part régulière de l'\omterm{prop:g3:balanced-diet:jobs}{énergie} et du dioxygène du corps. Privé
de l'un ou de l'autre, il réagit plus lentement, disperse son
attention et décide mal.
\end{solution}

\begin{solution}{exo:g5:brain-in-command:8}
Réponse libre. Les prises habituelles se situent autour de
$10$--$20$ cm~; plusieurs essais améliorent et stabilisent en
général un peu la valeur.
\end{solution}

\begin{solution}{exo:g5:brain-in-command:9}
Parce que toute prise doit parcourir les trois étapes --- le rapport
des \omterm{def:g1:the-five-senses:sense}{yeux} doit voyager, le \omterm{def:g5:brain-in-command:system}{cerveau} doit calculer et décider, l'ordre
doit revenir et les \omterm{def:g3:bones-and-muscles:muscle}{muscles} se contracter --- et que chaque étape
prend du temps réel. L'habileté raccourcit la boucle~; rien ne la
supprime.
\end{solution}

\begin{solution}{exo:g5:brain-in-command:10}
Parce que les meilleurs freins n'agissent qu'après que la boucle du
conducteur a tourné --- et qu'à vitesse d'autoroute la voiture
parcourt environ \qty{7}{m} avant même que l'ordre «~freine~!~»
n'atteigne le pied. L'espace entre les voitures est la place de la
boucle. L'attention raccourcit les étapes de rapport et de décision~;
un \omterm{def:g5:brain-in-command:system}{cerveau} distrait ajoute des mètres.
\end{solution}

\begin{solution}{exo:g5:brain-in-command:11}
La part du \omterm{def:g5:brain-in-command:system}{cerveau} --- la décision et la commande~: l'entraînement a
tracé de meilleurs chemins, si bien que les mêmes rapports
aboutissent maintenant vite et sans effort aux bons ordres.
L'entraînement d'un \omterm{def:g3:bones-and-muscles:muscle}{muscle} fait grandir la force du \omterm{def:g3:bones-and-muscles:muscle}{muscle}~;
l'entraînement d'un savoir-faire fait grandir l'aiguillage du
\omterm{def:g5:brain-in-command:system}{cerveau}. L'un bâtit le moteur, l'autre le chef.
\end{solution}
